\documentclass[11pt,a4paper]{article}

\usepackage[utf8]{inputenc}
\usepackage[T1]{fontenc}
\usepackage{lmodern}
\usepackage[margin=2.4cm]{geometry}
\usepackage{microtype}
\usepackage{hyperref}
\usepackage{xcolor}
\usepackage{booktabs}
\usepackage{longtable}
\usepackage{array}
\usepackage{enumitem}
\usepackage{listings}
\usepackage{fancyhdr}
\setlength{\headheight}{14pt}
\usepackage{titlesec}
\usepackage{tcolorbox}
\usepackage{amsmath}

\hypersetup{
  colorlinks=true,
  linkcolor=blue!55!black,
  urlcolor=blue!55!black,
  citecolor=blue!55!black,
  pdftitle={HMSpectralGun User Manual},
  pdfauthor={HMSpectralGun Team},
  pdfsubject={Operational manual for Turbospectrum-based stellar spectral synthesis and line-by-line fitting}
}

\definecolor{codebg}{RGB}{246,248,250}
\definecolor{rulegray}{RGB}{90,90,90}
\definecolor{warnbg}{RGB}{255,248,220}
\definecolor{notebg}{RGB}{239,247,255}

\lstdefinestyle{manualcode}{
  basicstyle=\ttfamily\small,
  backgroundcolor=\color{codebg},
  frame=single,
  rulecolor=\color{rulegray},
  breaklines=true,
  columns=fullflexible,
  keepspaces=true,
  showstringspaces=false,
  upquote=true
}
\lstset{style=manualcode}

\setlist[itemize]{topsep=3pt,itemsep=2pt,parsep=0pt,leftmargin=1.6em}
\setlist[enumerate]{topsep=3pt,itemsep=2pt,parsep=0pt,leftmargin=1.8em}
\renewcommand{\arraystretch}{1.2}

\pagestyle{fancy}
\fancyhf{}
\lhead{HMSpectralGun User Manual}
\rhead{Turbospectrum workflow}
\cfoot{\thepage}

\titleformat{\section}{\Large\bfseries}{\thesection}{0.7em}{}
\titleformat{\subsection}{\large\bfseries}{\thesubsection}{0.7em}{}
\titleformat{\subsubsection}{\normalsize\bfseries}{\thesubsubsection}{0.7em}{}

\newtcolorbox{importantbox}{
  colback=warnbg,
  colframe=orange!70!black,
  title=Important,
  fonttitle=\bfseries,
  boxrule=0.7pt,
  arc=2pt
}

\newtcolorbox{notebox}{
  colback=notebg,
  colframe=blue!55!black,
  title=Note,
  fonttitle=\bfseries,
  boxrule=0.7pt,
  arc=2pt
}

\newcommand{\code}[1]{\texttt{#1}}

\begin{document}

\begin{titlepage}
  \centering
  \vspace*{2.5cm}
  {\Huge\bfseries HMSpectralGun User Manual\par}
  \vspace{0.4cm}
  {\Large Turbospectrum-based stellar spectral synthesis and line-by-line abundance fitting\par}
  \vspace{1.5cm}
  {\large Improved LaTeX edition\par}
  \vfill
  \begin{tabular}{rl}
    Project: & \code{HMSpectralGun} \\
    Scope: & User-facing execution, input preparation, and validation workflow \\
    Status: & Operational manual \\
    Language: & English \\
  \end{tabular}
  \vfill
  {\large \today\par}
\end{titlepage}

\tableofcontents
\newpage

\section{Purpose and Scope}

This manual explains how to run the current \code{HMSpectralGun} workflows. The suite generates stellar synthetic spectra with Turbospectrum starting from atmospheric parameters, line lists, abundance patterns, and optional post-processing steps such as convolution, resampling, and noise injection.

The repository contains three main workflows:

\begin{longtable}{p{0.36\textwidth}p{0.54\textwidth}}
\toprule
\textbf{Workflow} & \textbf{Role} \\
\midrule
\code{main.py} & Classical serial synthesis workflow. It reads \code{input.ts} and synthesizes spectra one after another. \\
\code{main\_parallel.py} & Parallel synthesis workflow. It follows the same logic as \code{main.py}, but distributes independent jobs across multiple workers. \\
\code{line\_chi2\_on\_the\_fly.py} & Line-by-line abundance workflow. It builds trial grids in \code{[X/Fe]}, synthesizes spectra on the fly, and estimates the best fit using a chi-square criterion. \\
\bottomrule
\end{longtable}

\begin{importantbox}
Turbospectrum is an external dependency. \code{HMSpectralGun} prepares inputs, launch scripts, and post-processing products, but the physical synthesis is performed by the compiled Turbospectrum executables.
\end{importantbox}

\section{Recommended Quick Start}

From the repository root, first validate the Python environment and then run a minimal single-spectrum test before launching a long batch.

\begin{lstlisting}[language=bash]
python --version
python -c "import numpy, pandas, scipy, matplotlib; print('core stack: ok')"
python main.py --input input.ts --progress
\end{lstlisting}

For line-by-line analysis, start from the example JSON configuration and restrict the first run to one target:

\begin{lstlisting}[language=bash]
python line_chi2_on_the_fly.py --config line_fit_config.example.json --star <star_id>
\end{lstlisting}

\begin{notebox}
The first production rule is simple: never start a large batch before a one-spectrum or one-star validation run has produced a physically sensible output file and usable logs.
\end{notebox}

\section{Requirements}

\subsection{Python environment}

Python 3.10 or newer is recommended. The code uses the following Python packages:

\begin{itemize}
  \item \code{numpy}, \code{pandas}, \code{tqdm}, and \code{mendeleev};
  \item \code{scipy}, for resampling and optimization utilities;
  \item \code{matplotlib}, for line-by-line PDF summary reports;
  \item \code{PyAstronomy}, for instrumental convolution.
\end{itemize}

A minimal environment check is:

\begin{lstlisting}[language=bash]
which python
python - <<'PY'
import numpy, pandas, scipy, matplotlib
print('core scientific stack: ok')
PY
\end{lstlisting}

If instrumental convolution is required, also check:

\begin{lstlisting}[language=bash]
python - <<'PY'
import PyAstronomy
print('PyAstronomy: ok')
PY
\end{lstlisting}

\subsection{Turbospectrum installation}

Turbospectrum must be compiled and its executables must be reachable through the paths expected by the code.

\begin{importantbox}
Several paths are hardcoded in the current codebase, including paths such as \code{launchpath}, \code{dataset\_model\_path}, and the paths to \code{babsma\_lu} and \code{bsyn\_lu}. If the repository is moved to a different machine or directory, update these paths before running the workflows.
\end{importantbox}

\section{Input File Structure}

\subsection{The \code{input.ts} file}

The file \code{input.ts} is used by \code{main.py} and \code{main\_parallel.py}. The parser expects a rigid structure:

\begin{enumerate}
  \item line 1: \code{savepath};
  \item line 2: \code{linelistpath};
  \item line 3: \code{modelpath};
  \item line 4: \code{ExplicitModel=True|False};
  \item line 5: \code{interp=True|False|nearest};
  \item line 6: third keyword, named \code{NLTE} in the code;
  \item line 7 onward: spectrum table with either 15 or 17 columns.
\end{enumerate}

\begin{notebox}
According to the available operational notes, the \code{NLTE} keyword is read by the parser but is not applied inside \code{TurboSpecWriter}. Treat it as non-operational unless the active repository version proves otherwise.
\end{notebox}

Minimal template:

\begin{lstlisting}
/path/output_spectra/
/path/linelists/
/path/models/
ExplicitModel=False
interp=True
NLTE=False
3600,1.0  -1.00  0.20  15000  15500  2.0  st  *  28000  0.02  *  linelistH.ts  abu.ts  *  txt
\end{lstlisting}

\subsection{Spectrum table columns}

The spectrum table starts after the first six lines of \code{input.ts}. The expected column order is:

\begin{longtable}{p{0.08\textwidth}p{0.22\textwidth}p{0.60\textwidth}}
\toprule
\textbf{No.} & \textbf{Column} & \textbf{Meaning} \\
\midrule
\endhead
1 & \code{Model} & Explicit model filename, or \code{Teff,logg}. \\
2 & \code{[Fe/H]} & Model metallicity. \\
3 & \code{[a/Fe]} & Alpha enhancement. \\
4 & \code{lam\_i} & Initial wavelength. \\
5 & \code{lam\_f} & Final wavelength. \\
6 & \code{xi} & Microturbulent velocity. \\
7 & \code{chemistry} & Chemistry pattern: \code{st}, \code{ap}, \code{ae}, \code{mc}, or \code{hc}. \\
8 & \code{sampl} & \code{*} or resampling step in Angstrom. \\
9 & \code{RES} & Resolving power used for convolution. \\
10 & \code{resnum} & Numerical wavelength step used in the Turbospectrum synthesis. \\
11 & \code{monoelem} & \code{*}, atomic species such as \code{FeI}, or molecules such as \code{CO}, \code{OH}, \code{CN}. \\
12 & \code{linelist\_file} & File listing the line-list files to be used. \\
13 & \code{abu\_file} & Abundance file. \\
14 & \code{snr} & \code{*} or numerical signal-to-noise ratio. \\
15 & \code{extension} & Output extension, for example \code{txt}. \\
16 & \code{override\_elem} & Optional atomic number of the element to override. \\
17 & \code{override\_xfe} & Optional \code{[X/Fe]} value to force. \\
\bottomrule
\end{longtable}

If columns 16 and 17 are provided, the override has maximum priority over the abundance file.

\subsection{Abundance file: \code{abu.ts}}

The abundance file is a two-column, whitespace-separated table:

\begin{itemize}
  \item column 1: atomic number;
  \item column 2: abundance offset \code{[X/Fe]}.
\end{itemize}

Example:

\begin{lstlisting}
8   0.20
12  0.20
26  0.00
612613  15.0
\end{lstlisting}

The special row \code{612613} is interpreted as the isotopic ratio \code{C12/C13}.

\subsection{Line-list files}

The \code{linelist\_file} entry points to a plain-text file containing one line-list filename per line:

\begin{lstlisting}
vald_1.lin
vald_2.lin
\end{lstlisting}

In the classical synthesis workflows, these files are resolved relative to \code{linelistpath}. In the chi-square workflow, they are resolved relative to \code{linelist\_library\_path}.

\section{Workflow 1: Classical Serial Synthesis}

\subsection{Command}

Run the serial workflow with:

\begin{lstlisting}[language=bash]
python main.py --input input.ts
\end{lstlisting}

To display a progress bar:

\begin{lstlisting}[language=bash]
python main.py --input input.ts --progress
\end{lstlisting}

\subsection{Execution logic}

The serial workflow performs the following operations for each spectrum row:

\begin{enumerate}
  \item selects or prepares the atmospheric model, using interpolation, nearest-neighbor selection, or an explicit model;
  \item writes the Turbospectrum \code{.com} launch script;
  \item runs Turbospectrum;
  \item adds a descriptive header to the output spectrum;
  \item optionally applies convolution, resampling, and noise injection.
\end{enumerate}

Use this workflow when debugging a new configuration or when the number of spectra is small.

\section{Workflow 2: Parallel Synthesis}

\subsection{Command}

Run the parallel workflow with:

\begin{lstlisting}[language=bash]
python main_parallel.py --input input.ts
\end{lstlisting}

\subsection{Execution logic}

The parallel workflow uses \code{ProcessPoolExecutor} with up to \code{cpu\_count() - 1} workers. It is designed for independent spectra that can be synthesized concurrently.

Operational features include:

\begin{itemize}
  \item worker-level logging;
  \item collision avoidance by renaming final outputs with a \code{\_k<index>} suffix;
  \item a main log named \code{synth\_main\_YYYYMMDD\_HHMMSS.log}.
\end{itemize}

Use this workflow for large batches when the input file has already been validated in serial mode.

\section{Workflow 3: Line-by-line Chi-square Analysis}

\subsection{Configuration file}

The line-by-line workflow is controlled by a JSON configuration file. A typical starting point is:

\begin{lstlisting}
line_fit_config.example.json
\end{lstlisting}

Main configuration blocks include:

\begin{longtable}{p{0.28\textwidth}p{0.62\textwidth}}
\toprule
\textbf{Block} & \textbf{Role} \\
\midrule
\code{paths} & Observed spectra, synthesis directory, and results directory. \\
\code{keywords} & Global options, including the \code{[X/Fe]} trial grid, fitting windows, interpolation mode, and related settings. \\
\code{stars} & List of targets to analyze. \\
\bottomrule
\end{longtable}

\subsection{Command}

Basic execution:

\begin{lstlisting}[language=bash]
python line_chi2_on_the_fly.py --config line_fit_config.example.json
\end{lstlisting}

Useful runtime overrides:

\begin{lstlisting}[language=bash]
python line_chi2_on_the_fly.py --config cfg.json --star 301_4MCMC.H --xfe-min -0.5 --xfe-max 0.7 --xfe-step 0.1
\end{lstlisting}

\subsection{Typical outputs}

The usual outputs in \code{results\_path} are:

\begin{itemize}
  \item \code{<star>\_best.csv};
  \item \code{<star>\_grid.csv};
  \item \code{<star>\_summary.pdf}, if \code{matplotlib} is available;
  \item legacy files named \code{results\_<id>}.
\end{itemize}

\section{Output Naming and Generated Files}

Turbospectrum launch scripts are written as \code{<name>.com} files in \code{launchpath}. Final spectra are saved in \code{savepath} for the classical workflows, or under \code{analysis\_synth\_path/<star>/} for the chi-square workflow.

If \code{override\_elem} and \code{override\_xfe} are active, the generated filename includes an abundance tag such as \code{\_Fep010} or \code{\_Mgm025}.

\begin{notebox}
If a spectrum already exists and is non-empty, the workflow may skip its regeneration. This is intentional behavior and prevents accidental overwriting during repeated runs.
\end{notebox}

\section{Recommended Validation Flows}

Run tests in increasing order of cost.

\subsection{Test A: Environment import test}

\begin{lstlisting}[language=bash]
python -c "import numpy, pandas, scipy, matplotlib; print('ok')"
\end{lstlisting}

\subsection{Test B: One-spectrum serial synthesis}

Prepare an \code{input.ts} file with a single spectrum row, then run:

\begin{lstlisting}[language=bash]
python main.py --input input.ts --progress
\end{lstlisting}

Check that the output spectrum exists, is non-empty, and contains the expected header.

\subsection{Test C: Small parallel run}

After Test B passes, use a small multi-row input file:

\begin{lstlisting}[language=bash]
python main_parallel.py --input input.ts
\end{lstlisting}

Check the main log and verify that output filenames do not collide.

\subsection{Test D: One-star line-by-line run}

Run one target with a narrow trial grid:

\begin{lstlisting}[language=bash]
python line_chi2_on_the_fly.py --config cfg.json --star <star_id> --xfe-min -0.2 --xfe-max 0.2 --xfe-step 0.1
\end{lstlisting}

Inspect \code{<star>\_best.csv}, \code{<star>\_grid.csv}, and the summary PDF if generated.

\section{Troubleshooting}

\subsection{\code{ModuleNotFoundError}}

A Python dependency is missing from the active environment. Confirm the interpreter and install the missing package in that same environment.

\begin{lstlisting}[language=bash]
which python
python -c "import numpy, pandas, scipy, matplotlib; print('ok')"
\end{lstlisting}

\subsection{\code{Cannot parse MARCS model filename}}

The atmospheric model filename does not match the expected naming format. Use a valid MARCS filename or switch to a supported explicit model configuration.

\subsection{\code{No safe interpolation pair} or \code{No complete safe MARCS interpolation cube}}

The MARCS model dataset is insufficient for the requested atmospheric parameters. Add the missing models, change the requested parameters, or use nearest-neighbor selection if scientifically acceptable for the test.

\subsection{Missing or empty spectrum output}

Turbospectrum likely failed or did not write the expected output file. Check the generated \code{.com} file and the Turbospectrum log \code{.txt} files in \code{launchpath}.

\subsection{\code{Spectrum already exists} followed by skip}

The output file already exists and is non-empty. This is expected skip behavior. Remove the existing output only if regeneration is intended.

\section{Pre-run Checklist for Long Jobs}

Before launching a long synthesis or line-by-line run, verify the following:

\begin{itemize}
  \item absolute paths have been updated for the current machine;
  \item \code{contopac} is available or can be created;
  \item \code{input.ts} or the JSON configuration has been validated;
  \item line-list files resolve correctly;
  \item \code{abu.ts} resolves correctly and contains the intended abundances;
  \item Turbospectrum executables are reachable;
  \item a one-spectrum or one-star test has completed successfully;
  \item output directories are writable;
  \item old non-empty outputs have been removed if regeneration is required.
\end{itemize}

\section{Best Practices}

\begin{itemize}
  \item Use \code{main.py} first when debugging a new input file; move to \code{main\_parallel.py} only after the serial run works.
  \item Keep one minimal validation \code{input.ts} under version control for regression tests.
  \item Keep production outputs in dated directories to avoid mixing incompatible runs.
  \item Treat hardcoded absolute paths as configuration debt: document every required local edit before a run.
  \item For line-by-line analysis, start with a coarse \code{[X/Fe]} grid and refine only after confirming that the observed and synthetic spectra are aligned.
\end{itemize}

\appendix

\section{Minimal \code{input.ts} Template}

\begin{lstlisting}
/path/output_spectra/
/path/linelists/
/path/models/
ExplicitModel=False
interp=True
NLTE=False
3600,1.0  -1.00  0.20  15000  15500  2.0  st  *  28000  0.02  *  linelistH.ts  abu.ts  *  txt
\end{lstlisting}

\section{Minimal \code{abu.ts} Template}

\begin{lstlisting}
8   0.20
12  0.20
26  0.00
612613  15.0
\end{lstlisting}

\section{Operational Rule of Thumb}

For development and debugging, prefer the safest sequence:

\begin{lstlisting}[language=bash]
# 1. Validate imports
python -c "import numpy, pandas, scipy, matplotlib; print('ok')"

# 2. Run one serial spectrum
python main.py --input input.ts --progress

# 3. Only then run parallel or line-by-line workflows
python main_parallel.py --input input.ts
python line_chi2_on_the_fly.py --config cfg.json --star <star_id>
\end{lstlisting}

\end{document}
